\section{Removing all flag enhancements}\label{desc:section}

Work over $\mathbb C$, with the marked curve $(X;x_1,\ldots,x_s)$
and $U=X\setminus\{x_1,\ldots,x_s\}$.  For a parameter $\sigma$
on $U$, write $V_\sigma$ for its associated tensor local system.
We will pass from a nonzero
coherent eigencomplex on $\cA^+$ to a nonzero perverse eigenobject on
$\cA$.  The map
\[
                         q:\cA^+\longrightarrow\cA
\]
is a torsor under $H=T^s$.  Compactly supported direct image along
$q$ can kill a local system with nontrivial torus characters.  The
point of the argument is that unipotent boundary monodromy of the
eigenvalue forces unipotent monodromy along every factor of $H$.

We adapt the one-point descent argument of \cite[Section~7]{CT}.
Its local input is Gaitsgory's central-sheaf construction
\cite{GaitsgoryCentral} in the universal monodromic form of
Dhillon--Taylor \cite{DT}.  The extension needed here keeps the levels
at the other marked points \emph{transported} in every moving and
convolution correspondence.  We spell out the resulting natural
comparisons: the trace identity needs monodromy on one convolution
factor at a time.  The universal kernels occur only in Betti sheaves;
$q_!$ and the final perverse cohomology are operations on the original
geometric adic complexes.

\subsection{Monodromy on the full enhancement torus}

A complex on an $H$-torsor is \emph{monodromic along the torsor} if,
locally on its base, it is constructible for a product stratification
$\{S_\alpha\times H\}$ and, on a contractible patch of $H$, is pulled
back from a slice.  This includes the extension data between strata.
It allows arbitrary monodromy in $\pi_1(H)=X_*(T)^s$ and requires no
equivariant structure.  The condition is preserved by base change on
the base.

\begin{lemma}\label{desc:monodromic}
If $Q\in\Dlc(\cA^+,E)$ has nilpotent singular support, its Betti
realization is monodromic along the whole $T^s$-torsor $q$, and along
each enhancement separately.
\end{lemma}

\begin{proof}
By the cone description following Definition~\ref{geom:cone}, the
nilpotent cone on $\cA^+$ is the
smooth cotangent pullback of the ordinary-level cone: a nilpotent
Borel-valued residue has zero toral component at every marked point.
On a smooth torsor chart $S\times H$ the singular support is therefore
contained in $\Lambda_S\times0_H$.  The cone $\Lambda_S$ is isotropic
by Proposition~\ref{geom:dimension}.  Choose a microlocal
stratification of $S$ whose conormals contain it.  The product
stratification criterion used in \cite[proof of Proposition~6.3.2]{NY} and
\cite[Lemma~7.1]{CT} makes $Q$ constructible for
$\{S_\alpha\times H\}$.  Restriction to a slice over a contractible
patch of $H$ is an equivalence for this constructible category; hence
it identifies the entire complex with the pullback from that slice.
Continuation gives the assertion on overlaps.  These descriptions
persist under base change and, in particular, under passage to the
base of a single enhancement torus.
\end{proof}

\subsection{The single-disc inputs}

Fix one marked point $x_i$.  In this subsection $T$ denotes its
individual enhancement torus; put
\[
 d=\dim T,\qquad \Gamma=X_*(T),\qquad
 R_T=E[\Gamma]=\mathcal O(\widehat T).
\]
Here $\widehat T\subset\Gd$ is the dual maximal torus.
We fix positive cocharacter loops and orientations of the compact
real torus in $T(\mathbb C)$.  Let $\mathscr L_T$ be the local system
with fiber $R_T$ and regular monodromy.  It has infinite-dimensional
stalks and is used in the category of weakly constructible Betti
sheaves.  For ordinary-direct-image convolution its unit is
\begin{equation}\label{desc:unit}
                         \mathbf1_T=\mathscr L_T[d].
\end{equation}
Our enhancement-difference convention writes the varying input as
$e a^{-1}$ for fixed output $e$ and difference $a\in T$.

We use the following precise features of this unit, proved in
\cite[Lemma~7.2]{CT}.  Convolution with $\mathbf1_T$ is naturally the
identity on relatively monodromic complexes, commutes with base
change, and respects successive convolutions.  Multiplication by
$R_T$ on the unit acts as enhancement monodromy.  In fact, if $W$ is
a fiber and $M_1,\ldots,M_d$ its commuting monodromies, the unit
integral before shifting is the Koszul complex
\[
 K^\bullet(z_1M_1^{-1}-1,\ldots,z_dM_d^{-1}-1;
                                      R_T\otimes_E W).
\]
Its cohomology is $W$ in degree $d$, with residual $z_j$ acting by
$M_j$.  This also applies to complexes relatively over a stratified
base.  Compatibility with two integrations follows on the real
universal covers of the compact tori: replace two lifted differences
by the first and their sum, and integrate in that order.  The change
preserves product orientation and continuation follows the same
summed path.  This specifies the unit maps, rather than only their
isomorphism classes.

Let $\mathrm{Iw}$ and $\mathrm{Iw}^0$ be the inverse images of $B$ and
$N_B$ in $L^+G$.  The local central-sheaf input is the monoidal functor
\[
 Z:\Rep(\Gd)\longrightarrow\mathcal H^{\mathrm{mon}},\qquad V\longmapsto Z(V),
\]
where $\mathcal H^{\mathrm{mon}}$ is the universal monodromic Betti
Hecke category on $LG/\mathrm{Iw}^0$, with the pro-unipotent
equivariance convention of \cite{DT}.  It comes with nearby-cycle
monodromies $m_V$, natural in $V$.  We use the kernel specialization
and monoidal comparison of \cite[Proposition~6.3.1 and
Lemma~6.3.2]{DT}, as recorded in \cite[Sections~7.2 and~7.4]{CT}:
the moving Satake kernel tensored with $\mathbf1_T$ specializes to
$Z(V)$, and the specialized twisted product and its convolution map
are the specified monoidal maps of $Z$.

The corresponding local trace formula is
\begin{equation}\label{desc:local-trace}
 \left(\mathbf1_T\xrightarrow{\mathrm{coev}}Z(V)\star Z(V^*)
 \xrightarrow{m_V\star\mathrm{id}}Z(V)\star Z(V^*)
 \xrightarrow{\mathrm{ev}}\mathbf1_T\right)
                 =\chi_V\in\operatorname{End}(\mathbf1_T)=R_T,
\end{equation}
where $\chi_V$ is the character on $\widehat T$ and evaluation in this
order is induced from the symmetric representation category.
For completeness, \cite[Propositions~5.3.1 and~9.3.1(a)--(c)]{DT}
compute this trace after the Wakimoto associated-graded functor.
On the weight summand $\nu\in X^*(\widehat T)=\Gamma$, monodromy is
multiplication by $e^\nu$.  The trace is thus
$\sum_\nu\dim(V_\nu)e^\nu=\chi_V$.  On endomorphisms of the unit the
associated-graded functor retains multiplication on its regular
$R_T$-module, and is faithful there, proving
\eqref{desc:local-trace}.  Reversing conventions inverts both the
torus and boundary loops and does not change the unipotence conclusion.

One further input explains why ordinary pushforward can be
specialized in these correspondences.  We state it because its
hypothesis must be checked again when the other levels are present.

\begin{lemma}[Exchange through an enhancement torus]\label{desc:exchange}
Let $P\xrightarrow{\pi}C\xrightarrow{b}Y$ be a factorization over an
analytic disc $\Delta$, with $\pi$ a torus torsor on the entire family
and $b$ proper.  On finite-dimensional paracompact charts suppose that
$F$ over $\Delta^*$ is weakly constructible and relatively monodromic
along $\pi$.  Then the natural exchange
\begin{equation}\label{desc:push-exchange}
            \Psi(b\pi)_*F\longrightarrow(b_0\pi_0)_*\Psi F
\end{equation}
is an isomorphism.  Base change and projection formula for these
pushforwards also hold, and the comparisons respect composition.
Here $\Psi$ denotes geometric Betti nearby cycles, without taking
monodromy invariants.
\end{lemma}

\begin{proof}
This is \cite[Lemma~7.3]{CT}.  Quotient the enhancement torus by its
positive real radial subgroup $A\simeq\mathbb R^d$, to factor
$P\xrightarrow{\rho}\overline P\xrightarrow{\bar\pi}C$ over the
\emph{whole} disc.  The second map is a compact-torus bundle and
therefore proper.  Relative monodromicity gives
$F=\rho^*\overline F$: restriction across a contractible radial
factor is an equivalence, so the local descriptions glue.  In
trivializations extending across zero,
\[
 R\rho_*\rho^*\overline F=\overline F,
            \qquad \Psi F=\rho_0^*\Psi\overline F.
\]
These identities persist after base change and tensoring with a
factor from the target.  Proper base change for $b\bar\pi$ proves
all assertions.  The maps are the natural exchanges for restriction
and direct image from the universal cover of $\Delta^*$, so they
respect composition.  The same proof applies to products of tori.
\end{proof}

\subsection{A central action with all the other levels retained}

Choose a disc $\Delta$ about $x_i$, disjoint from the other marks,
and a lift to the universal cover of $\Delta^*$.  For the eigenvalue
$V_\sigma$, write $V_{\partial i}$ for the resulting boundary fiber
and $M_{V,i}$ for positive-loop monodromy.  Let $\mu_i$ denote the
$R_T$-action given by continuation along the $i$th enhancement.

\begin{lemma}[Trace at one of several marked points]\label{desc:trace}
Let $Q$ be a locally bounded algebraically constructible Betti
complex on $\cA^+$, monodromic along $T^s$,
with coherent Hecke eigenisomorphisms for $\sigma$ on $U$.  The local
central kernels at $x_i$ act on $Q$, and there are natural
isomorphisms
\begin{equation}\label{desc:central-eigen}
          Z(V)\star_{x_i}Q\simeq Q\otimes_E V_{\partial i}
\end{equation}
compatible with the unit, representation morphisms, and convolution.
On each individual convolution factor they carry $m_V$ to $M_{V,i}$.
In particular, for every $i$ and every $V\in\Rep(\Gd)$,
\begin{equation}\label{desc:character-trace}
              \chi_V(\mu_i)=\operatorname{Tr}(M_{V,i})\,\mathrm{id}_Q.
\end{equation}
\end{lemma}

\begin{proof}
We give the several-level version of \cite[Lemma~7.4]{CT}, including
the comparisons needed to apply \eqref{desc:local-trace}.

\emph{Moving correspondence and its single-step comparison.}
Let $\mathfrak H_{i,\Delta}$ parametrize a leg $z\in\Delta$, two
bundles $P_0,P_1$ with enhancements $\eta_{0j},\eta_{1j}$ at every
$x_j$, and a modification $P_0\simeq P_1$ off $z$.  At $x_i$ the two
enhancements are independent.  At each $x_j$ with $j\ne i$ require
$\eta_{1j}$ to be the transport of $\eta_{0j}$.  Write $p$ for input
and $o$ for output-and-leg.  We bound the underlying modifications
by closed spherical Schubert bounds, retaining the full independent
flags at $x_i$.  Choose these bounds to contain the Satake supports;
their closed central fibers then contain the specialized supports.

Forgetting the levels at $j\ne i$ makes this the pullback of the
one-point moving correspondence by the stack of those levels on
its input.  Transport gives the same description using its output.
This cartesian assertion also holds for successive modifications:
there is one choice of each other enhancement, and the intermediate
and final choices are determined.  It does \emph{not} assert that
$Q$ is pulled back from the one-point stack.

In an input frame compatible with $\eta_{0i}$ the local family is
the one used to define $Z$: its punctured fiber is
$\Gr_G\times G/N_B$ and its central fiber is $LG/\mathrm{Iw}^0$.
The kernel $K_V$ on the punctured fiber is the relative Satake
kernel on $\Gr_G$, tensored with $\mathbf1_T$ on
$T=B/N_B\subset G/N_B$, extended by zero.  Thus it requires equality
of the two ordinary flags at $x_i$ and records their enhancement
difference.  Its specialization is $Z(V)$.  With the relative
normalization of Section~\ref{found:section}, the moving-leg
perverse shift is absent on both sides of this statement.

The kernels and their maps descend from these frames.  On a bound
one can take finite jets on a sufficiently thick neighborhood of
the sum of the graphs of $x_i$ and $z$, so the same smooth frame
space works at collision.  The equivariance of Satake and the
$N_B$-invariance of the flag factor give descent off zero; at zero
it is the $\mathrm{Iw}^0$-equivariance of $Z(V)$.  The finite jet
quotients of this last group are unipotent, hence Betti contractible,
so descent retains the maps and their cocycles.  The extra levels
are outside this neighborhood and enter only in the bundle factor.

For an input family $D$ over $\Delta^*$, write
\[
 A_V(D)=o_*(D\,\widetilde\boxtimes K_V),
                 \qquad B_V(R)=Z(V)\star_{x_i}R,
\]
where the twisted product includes the fixed relative normalization.
For $Q$ viewed as a constant family, the unit calculation gives
\begin{equation}\label{desc:punctured-action}
 A_V(Q)\simeq\Hecke_V(Q)|_{\Delta^*}
                    \simeq Q\boxtimes V_\sigma|_{\Delta^*}.
\end{equation}
Indeed the only integration beyond the usual Hecke correspondence
is over the input enhancement at $x_i$, with kernel $\mathbf1_T$.

The output map $o$ factors over the full disc by forgetting just
$\eta_{0i}$, keeping its ordinary flag, and then by a proper map.
The first map is a $T$-torsor.  For the second, bounded input
modifications with fixed output lie in a proper Grassmannian bound,
and their ordinary input flags form a proper $G/B$-bundle.  All
other enhancements are determined by transport from the output;
they add no nonproper directions.  At $z=0$ the output enhancement
is still fixed, and choosing the ordinary input flag is still the
full proper flag bundle on the modified input fiber.  Thus the
quotient is proper at collision as well; we have not replaced it
by an exact-position affine-flag stratum.
The integrand is relatively monodromic along the forgotten torus:
$Q$ has this property by hypothesis, and $K_V$ has a local system in
the enhancement difference with support conditions independent of
that difference.  Lemma~\ref{desc:exchange} therefore applies.

In the input frames, the nearby-cycle tensor map for this integrand
is an exterior-product comparison with an arbitrary coefficient
from the bundle factor.  It is an isomorphism both for constant
$D=Q$ and for $D=Q\boxtimes L$, where $L$ is a finite-rank local
system on $\Delta^*$: trivialize $L$ on the puncture cover.  This is
a Betti calculation on finite-dimensional bounds, as in
\cite[Lemma~7.4]{CT}.  On sufficiently small chart neighborhoods,
the local Milnor data
adapted to the constructible factors have finite triangulations.
Finite cellular complexes then compute the comparison, permitting
the infinite-dimensional stalks of $\mathscr L_T$; no global finite
triangulation of the correspondence is needed.  There is no adic
inverse limit in this calculation.  Smooth frame descent preserves
the comparison and its naturality.

Combining this tensor map with inverse push exchange defines an
actual natural isomorphism
\begin{equation}\label{desc:comparison}
 \begin{split}
 c_V(D):B_V(\Psi D)
 &=o_{0,*}(\Psi D\,\widetilde\boxtimes Z(V))\\
 &\longrightarrow o_{0,*}\Psi(D\,\widetilde\boxtimes K_V)
       \xrightarrow{\sim}\Psi A_V(D).
 \end{split}
\end{equation}
For constant $D=Q$, its composite with the specialization of
\eqref{desc:punctured-action} gives \eqref{desc:central-eigen} and
identifies $m_V$ with $M_{V,i}$.  We still need to check that this
single-step identification remains valid on each factor of a
convolution.

\emph{The two-step comparison.}
First, for $D=Q\boxtimes L$ the construction gives the commutative
square
\begin{equation}\label{desc:lisse-square}
\begin{CD}
 B_V(\Psi(Q\boxtimes L)) @>{c_V(Q\boxtimes L)}>>
                         \Psi A_V(Q\boxtimes L)\\
 @V{\sim}VV @VV{\sim}V\\
 B_V(Q)\otimes L_\partial @>{c_V(Q)\otimes\mathrm{id}}>>
                         \Psi A_V(Q)\otimes L_\partial.
\end{CD}
\end{equation}
On the puncture cover both routes are the same tensor comparison
for $Q,K_V$, tensored with the fiber of $L$, followed by the same
push exchange.  Finite rank allows that fiber through the direct
images.  The square is natural in every local-system endomorphism
of $L$, including its monodromy: the latter commutes with the cyclic
fundamental group of $\Delta^*$ and is an endomorphism of $L$ itself.

Now let $\mathfrak H^{(2)}_{i;W,V}$ parametrize successive bounded
modifications $P_0\to P_1\to P_2$ at the same $z\in\Delta$.
The three enhancements at $x_i$ are independent; those at every
other mark are transported throughout.  Its integrand is
$Q\,\widetilde\boxtimes K_W\,\widetilde\boxtimes K_V$.
There are two useful pushes.  The map $f$ forgetting $P_0$ and the
first modification is the base change of the first output map.
It factors by forgetting $\eta_{0i}$, then by a proper map.  The map
$g$ composing the modifications and forgetting $P_1$ factors by
forgetting $\eta_{1i}$, then by a proper map: the bounded intermediate
bundles form the closed convolution fiber in a proper Grassmannian
bound, and the intermediate ordinary flag contributes a proper
$G/B$-bundle.  Choose a target bound containing all composites.
Both factorizations are over the whole disc.  Transport again
determines the other marked-point enhancements.

The integrand is relatively monodromic for both forgotten tori.
For $f$ this is the single-step check, with the second kernel pulled
back from the target.  For $g$, $Q$ is independent of $\eta_{1i}$,
and varying $\eta_{1i}$ varies the two enhancement differences
inversely.  Their universal local systems remain a local system in
this variable; their ordinary-flag and Schubert conditions are
independent of it.  Lemma~\ref{desc:exchange} gives the push,
base-change, and projection-formula comparisons for $f$ and $g$.

The unit convolution on $\Delta^*$ and the specified monoidal map
of $Z$ at zero induce
\[
 b:A_VA_W(Q)\xrightarrow{\sim}A_{V\otimes W}(Q),\qquad
 b_0:B_VB_W(Q)\xrightarrow{\sim}B_{V\otimes W}(Q).
\]
Tensor factors here are in action order.  The essential compatibility
is
\begin{equation}\label{desc:two-step-square}
\begin{CD}
 B_VB_W(Q) @>{b_0}>> B_{V\otimes W}(Q)\\
 @V{c_V(A_W(Q))\,B_V(c_W(Q))}VV
                 @VV{c_{V\otimes W}(Q)}V\\
 \Psi(A_VA_W(Q)) @>{\Psi b}>>\Psi A_{V\otimes W}(Q).
\end{CD}
\end{equation}
The left side is defined because \eqref{desc:punctured-action}
identifies $A_W(Q)$ with $Q\boxtimes W_\sigma|_{\Delta^*}$, an
allowed input for $c_V$.

To prove the square, start on the two-step correspondence with the
natural map from the twisted product of the nearby cycles of
$Q,K_W,K_V$ to the nearby cycles of their twisted product.
The local kernel comparison of \cite[Lemma~6.3.2]{DT}, together
with exterior comparison for $Q$, makes it an isomorphism.
Pushing first through $f$ identifies it with the left side of
\eqref{desc:two-step-square}, by base change and projection formula.
Pushing first through $g$ identifies it with $b_0$ followed by
$c_{V\otimes W}(Q)$: the specialized kernel map is precisely the
monoidal map of $Z$ in the input frames.  After the final output
push these are equal, since tensor exchange, push exchange, and
the intervening projection-formula squares are natural and compose.
Every nonproper exchange in this assertion has one of the torus
factorizations just checked.  On the punctured family the same
equality uses the compatible successive unit integrations specified
after \eqref{desc:unit}.  Thus the square compares these particular
maps, not merely two isomorphic outputs.

\emph{Individual monodromy and trace.}
Apply \eqref{desc:lisse-square} with $L=W_\sigma|_{\Delta^*}$.
The comparison $c_W(Q)$ sends $m_W$ to $M_{W,i}$; naturality of the
square in that local-system endomorphism carries it to the
$W_{\partial i}$ factor alone after applying $B_V$.
For the outer factor, the left vertical identification in
\eqref{desc:lisse-square} is a projection-formula map, natural in
endomorphisms of the kernel $Z(V)$.  It therefore intertwines
$m_V$ acting on $B_V(\Psi(Q\boxtimes L))$ with $m_V$ on $B_V(Q)$
tensored with $\mathrm{id}_{L_\partial}$.  The bottom arrow is
$c_V(Q)\otimes\mathrm{id}$, which identifies this operator with
$M_{V,i}$ on its factor alone.  More explicitly, let
\[
 \Phi_{V,W}:B_VB_W(Q)\xrightarrow{\sim}
                         Q\otimes V_{\partial i}\otimes W_{\partial i}
\]
be the successive comparison and eigenidentification, with the
fixed symmetry placing the factors in action order.  What we have
proved is the pair of identities
\begin{align*}
 \Phi_{V,W}(m_V\star\mathrm{id})
   &=(\mathrm{id}_Q\otimes M_{V,i}\otimes\mathrm{id})\Phi_{V,W},\\
 \Phi_{V,W}(\mathrm{id}\star m_W)
   &=(\mathrm{id}_Q\otimes\mathrm{id}\otimes M_{W,i})\Phi_{V,W}.
\end{align*}
The square
\eqref{desc:two-step-square} and the original coherent eigenmaps
identify this comparison with the one for $V\otimes W$.
Naturality in representation morphisms includes coevaluation and
evaluation; the tensor unit is covered by \eqref{desc:unit}.
Iteration gives the same assertion for any number of factors.

Consequently, applying \eqref{desc:local-trace} to $Q$ gives insertion
and contraction of $V_{\partial i}$ with $M_{V,i}$ on its indicated
factor, hence $\operatorname{Tr}(M_{V,i})\mathrm{id}_Q$.
The right side of that local formula acts by $\chi_V(\mu_i)$ through
the universal-unit identification.  This proves
\eqref{desc:character-trace} with all other enhancements retained.
\end{proof}

\subsection{Unipotence, nonvanishing, and the ordinary-level object}

\begin{proposition}\label{desc:ordinary}
Let $\sigma$ be a continuous Zariski-dense geometric adic parameter
on $U$, defined over a finite coefficient extension, with unipotent
local monodromy at every $x_i$.  If $\cA^+$ carries a nonzero lcc
geometric adic coherent Hecke eigencomplex for $\sigma$, then $\cA$
carries a nonzero locally constructible perverse geometric adic
Hecke eigensheaf for $\sigma$.  It has all the unit, convolution,
permutation, and fusion compatibilities of Definition~\ref{found:coherence},
and nilpotent singular support.  No boundary monodromy is required
to be regular.
\end{proposition}

\begin{proof}
We extend the final step of \cite[Proposition~7.5]{CT} to the full
$T^s$-torsor.  Take a nonzero perverse cohomology $Q$ of the given
eigencomplex.
Proposition~\ref{perv:eigen} supplies its coherent eigenstructure,
and Theorem~\ref{det:field} puts its singular support in the
nilpotent cone.  Lemma~\ref{desc:monodromic} and Betti comparison
therefore allow us to apply Lemma~\ref{desc:trace} at every $x_i$.
Since $M_{V,i}$ is unipotent,
\begin{equation}\label{desc:character-dimension}
                       \chi_V(\mu_i)=(\dim V)\mathrm{id}_Q
                     \qquad\text{for every }i,V.
\end{equation}

Restrict to one $H=T^s$ fiber and to a finite-dimensional stalk
cohomology space $W$ there.  Its commuting monodromies have joint
generalized eigencharacters, each a tuple
$(a_1,\ldots,a_s)\in\widehat T(E)^s$.  Equation
\eqref{desc:character-dimension} says $\chi_V(a_i)=\chi_V(1)$ for
all representations $V$.  Characters span
$\mathcal O(\widehat T)^{W_G}$, where $W_G$ is the Weyl group,
and separate semisimple classes.
Thus $a_i$ has the same image as $1$ in the finite Weyl quotient;
the fiber through $1$ is the singleton $\{1\}$.  Every $a_i=1$,
so all commuting monodromy operators along the whole $H$ are
unipotent.  This conclusion uses no regularity or fixed Jordan type.

Return to adic sheaves and put $F=q_!Q$.  The representable morphism
$q$ is of finite type with fixed relative dimension, so $F$ is lcc.
To prove $F\ne0$, choose a fiber $H=q^{-1}(a)$ on which $Q$ has a
nonzero stalk.  Compute on its Betti realization, whose cohomology
sheaves are finite-rank local systems.
Let $j$ be the largest index with $\mathcal H^j(Q|_H)\ne0$, and
let $W$ be its fiber.  Write $r=\dim H=s\dim T$ and
$\Gamma_H=\pi_1(H^{\mathrm{an}})$.  Then
\[
                  H_c^{2r}(H,\mathcal H^j(Q|_H))
                              \simeq W_{\Gamma_H}(-r)\ne0.
\]
Indeed the finitely many commuting operators $M-1$ are nilpotent,
so the augmentation ideal $\mathfrak a\subset E[\Gamma_H]$ acts
nilpotently on this particular $W$.  If its coinvariants vanished,
$W=\mathfrak aW=\mathfrak a^2W=\cdots=0$, a
contradiction.  The displayed identification is Poincar\'e duality.
In the compact-support hypercohomology spectral sequence
\[
 E_2^{a,b}=H_c^a(H,\mathcal H^b(Q|_H))
                              \Longrightarrow H_c^{a+b}(H,Q|_H),
\]
the nonzero term $(a,b)=(2r,j)$ has no incoming differential, because
higher cohomology sheaves vanish, and no outgoing differential,
because $H_c^a$ vanishes for $a>2r$.  It survives, proving nonzero
compactly supported hypercohomology.  Betti comparison and base
change for $q_!$ identify this with a nonzero stalk of $F$.

Finally the unramified Hecke correspondences transport all levels.
The enhanced correspondence is cartesian from the ordinary one by
$q$ using either projection, and has the same relative Satake
kernel.  Base change and projection formula give, for every leg set
$I$,
\begin{align*}
 \Hecke_{I,(V_i)}(q_!Q)
 &\simeq(q\times\mathrm{id}_{U^I})_!\Hecke_{I,(V_i)}(Q)\\
 &\simeq q_!Q\boxtimes\mathop{\boxtimes}_{i\in I}(V_i)_\sigma.
\end{align*}
On bounds the Hecke pushforwards are proper.  The same cartesian
description holds on successive-modification diagrams and after
restriction to every collision diagonal; their natural exchanges
compose.  Thus these isomorphisms transport all the coherence
diagrams, in the original relative normalization without a
moving-leg shift.  The nonzero lcc eigencomplex $F$ has a nonzero
perverse cohomology.  Apply Proposition~\ref{perv:eigen} to obtain
the asserted coherent eigensheaf and Theorem~\ref{det:field} for
its nilpotent singular support.
\end{proof}
